\documentclass[12pt]{article}
\usepackage[utf8]{inputenc}\usepackage[T1]{fontenc}
\usepackage[a4paper,margin=2.5cm]{geometry}
\usepackage{amsmath,amssymb,amsthm,mathtools,mathrsfs}
\usepackage[dvipsnames]{xcolor}
\usepackage{booktabs,enumitem,hyperref}
\setlist{nosep,leftmargin=1.8em}
\hypersetup{colorlinks=true,linkcolor=blue!60!black,citecolor=green!40!black,urlcolor=blue!50!black}

\theoremstyle{plain}
\newtheorem{theorem}{Theorem}[section]
\newtheorem{proposition}[theorem]{Proposition}
\newtheorem{corollary}[theorem]{Corollary}
\newtheorem{lemma}[theorem]{Lemma}
\theoremstyle{definition}
\newtheorem{definition}[theorem]{Definition}
\newtheorem{axiom}[theorem]{Axiom}
\theoremstyle{remark}
\newtheorem{remark}[theorem]{Remark}

\newcommand{\Poinc}{\mathrm{Poinc}}
\newcommand{\SO}{\mathrm{SO}}
\newcommand{\SU}{\mathrm{SU}}
\newcommand{\SL}{\mathrm{SL}}
\newcommand{\ISO}{\mathrm{ISO}}
\newcommand{\Diff}{\mathrm{Diff}}
\newcommand{\Ker}{\mathrm{Ker}}
\newcommand{\Stab}{\mathrm{Stab}}
\newcommand{\Ad}{\mathrm{Ad}}
\newcommand{\Spec}{\mathrm{Spec}}
\newcommand{\Tr}{\mathrm{Tr}}
\newcommand{\dd}{\mathrm{d}}
\newcommand{\hhat}{\hat}

\title{\textbf{Quantum Mechanics from Orbital Structure:\\The Schr\"odinger and Dirac Equations from\\Unitarity, Representations, and Conservation}}

\author{Jocelyn Nemb\'e\thanks{Email: jnembe777@gmail.com, jnembe@root-insights.org}\\[0.3em]
\small Roots Insights Laboratory, Singapore\\
\small National Institute of Management Sciences, Libreville, Gabon}

\date{}

\begin{document}
\maketitle

% ============================================================================
\begin{abstract}
We derive the fundamental wave equations of quantum mechanics --- the Schr\"odinger, Klein--Gordon, and Dirac equations --- from the same three-step orbital pattern used to derive special relativity (Paper~1) and general relativity (Paper~2): \emph{symmetry group $\to$ conservation $\to$ uniqueness}. In the quantum regime ($\beta$), the conservation principle becomes \textbf{unitarity} (probability conservation), the symmetry group is the \textbf{Bargmann group} (central extension of the Galilei group) for non-relativistic mechanics or the \textbf{Poincar\'e group} for relativistic mechanics, and the uniqueness theorem is \textbf{Schur's lemma} applied to the Casimir operators.

We prove three results: (i)~the \textbf{unitarity--Noether correspondence} (Theorem~\ref{thm:correspondence}): unitarity in quantum mechanics is the precise structural analog of Noether conservation in classical physics, connected by Stone's theorem; (ii)~the Schr\"odinger equation is the \emph{unique} evolution equation compatible with Bargmann equivariance, unitarity, and irreducibility, via the Casimir $C_2 = H - P^2/2m$ and Schur's lemma (Theorem~\ref{thm:schrodinger}); (iii)~the \textbf{Bianchi--Casimir correspondence} (Theorem~\ref{thm:bianchi-casimir}): the Bianchi identity $\nabla_\mu G^{\mu\nu} \equiv 0$ (which determines Einstein's equations in Paper~2) and the Casimir eigenvalue equation $C_i = \lambda_i \cdot \mathrm{Id}$ (which determines the wave equations) play structurally identical roles as \emph{orbital integrability conditions}.

No Lagrangian is used. The theories are classified by a pair $(\kappa, \rho)$: kernel dimension $\kappa$ and coupling regime $\rho \in \{\alpha, \beta\}$. SR $= (10, \alpha)$; GR $= (<\!10, \alpha)$; QM $= (10, \beta)$; quantum gravity $= (<\!10, \beta)$.
\end{abstract}

\noindent\textbf{Keywords:} Schr\"odinger equation, Bargmann group, Wigner classification, unitarity, Noether conservation, orbital structure, Casimir operators, Stone--von Neumann, Lagrangian-free derivation

\noindent\textbf{MSC 2020:} 81R05, 22E70, 81Q70


% ============================================================================
\section{Introduction}
\label{sec:intro}
% ============================================================================

\subsection{The standard derivations}

The Schr\"odinger equation was introduced heuristically~\cite{schrodinger1926} via the de Broglie--Hamilton--Jacobi analogy: replace the classical energy $E$ by $i\hbar\partial_t$ and momentum $p$ by $-i\hbar\nabla$ in $E = p^2/2m + V$. The Dirac equation~\cite{dirac1928} was obtained by factorizing the Klein--Gordon operator. Feynman's path integral~\cite{feynman1948} derives the Schr\"odinger equation from a Lagrangian via the sum over paths. All three approaches involve either heuristic substitution rules or variational principles.

A different tradition derives wave equations from group representations. Wigner~\cite{wigner1939} classified elementary particles as irreducible unitary representations of the Poincar\'e group. Bargmann~\cite{bargmann1954} identified the central extension of the Galilei group as the correct symmetry for non-relativistic quantum mechanics, with mass as the central charge. L\'evy-Leblond~\cite{levyleblond1967} constructed Galilei-invariant wave equations for arbitrary spin. Musielak~\cite{musielak2021} derived infinite sets of wave equations from the extended Galilean group.

\subsection{The question}

Papers~1 and 2~\cite{nembe2026p1,nembe2026p2} derived SR and GR from the orbital pattern: symmetry $\to$ conservation $\to$ uniqueness, without variational principles. Does the \emph{same} pattern produce quantum mechanics? Specifically:
\begin{itemize}
    \item What plays the role of ``conservation'' in the quantum regime?
    \item What plays the role of ``uniqueness''?
    \item Is the structural parallel with GR precise or merely analogical?
\end{itemize}

\subsection{The answers}

Conservation becomes \textbf{unitarity}: $\|U(t)\psi\| = \|\psi\|$ for all $t$ (probability is conserved). Uniqueness becomes \textbf{Schur's lemma}: on an irreducible representation, the Casimir operators take definite scalar values, which determine the Hamiltonian and hence the wave equation.

The parallel with GR is precise, not analogical. The Bianchi identity (Paper~2) and the Casimir eigenvalue equation (this paper) play the same structural role: both are \emph{orbital integrability conditions} that constrain the dynamics from the group structure alone.


% ============================================================================
\section{The Regime Transition: Classical to Quantum}
\label{sec:regime}
% ============================================================================

\begin{definition}[Two regimes of equivariance]\label{def:regimes}
Let $G$ be a symmetry group acting on a state space $\mathcal{S}$.
\begin{itemize}
    \item \textbf{Regime $(\alpha)$} (classical): The state is a point $s \in \mathcal{S}$. Observables are functions $f: \mathcal{S} \to \mathbb{R}$. The group acts by diffeomorphisms: $g \cdot s \in \mathcal{S}$. The coupling between symmetry and observation is \emph{deterministic}: a transformation $g$ produces a definite change $f(g \cdot s)$.
    \item \textbf{Regime $(\beta)$} (quantum): The state is a ray $[\psi] \in \mathbb{P}\mathcal{H}$ in a projective Hilbert space. Observables are self-adjoint operators $\hat{A}$ on $\mathcal{H}$. The group acts by unitary operators: $g \mapsto U(g)$. The coupling is \emph{stochastic}: a measurement of $\hat{A}$ on state $\psi$ yields eigenvalue $a_i$ with probability $|\langle a_i|\psi\rangle|^2$.
\end{itemize}
\end{definition}

\begin{remark}[What changes between regimes]
The \emph{group} $G$ is the same (Poincar\'e or Galilei). What changes is the \emph{representation space} and the \emph{coupling type}:
\begin{center}\small
\begin{tabular}{lcc}\toprule
& \textbf{Regime $(\alpha)$} & \textbf{Regime $(\beta)$} \\\midrule
State & Point $x \in \mathcal{M}$ & Vector $\psi \in \mathcal{H}$ \\
Observable & Function $f(x)$ & Operator $\hat{A}$ \\
Symmetry action & Diffeomorphism $\psi^*$ & Unitary $U(g)$ \\
Measurement & Deterministic $f(x)$ & Stochastic $|\langle a_i|\psi\rangle|^2$ \\
\bottomrule
\end{tabular}
\end{center}
The structural constant of regime $(\beta)$ is $\hbar$: it parameterizes the central extension (Bargmann) and sets the scale of non-commutativity. Like $c$ (Paper~1) and $G, \Lambda$ (Paper~2), $\hbar$ is an \emph{input parameter}, not derived from the framework.
\end{remark}


% ============================================================================
\section{Unitarity as Quantum Conservation}
\label{sec:unitarity}
% ============================================================================

\begin{theorem}[Stone's theorem~\cite{stone1932,reed1972}]\label{thm:stone}
Every strongly continuous one-parameter unitary group $\{U(t)\}_{t \in \mathbb{R}}$ on a Hilbert space $\mathcal{H}$ has the form:
\begin{equation}
U(t) = e^{-iHt/\hbar}
\end{equation}
where $H$ is a (possibly unbounded) self-adjoint operator on $\mathcal{H}$, called the \emph{generator}. Conversely, every self-adjoint operator generates a strongly continuous unitary group.
\end{theorem}

\begin{theorem}[Unitarity--Noether correspondence]\label{thm:correspondence}
The structural analog between classical conservation and quantum unitarity is:
\begin{center}\small
\begin{tabular}{lll}\toprule
\textbf{Concept} & \textbf{Classical (regime $\alpha$)} & \textbf{Quantum (regime $\beta$)} \\\midrule
Structure preserved & Metric $g$ (by isometries) & Norm $\|\psi\|$ (by unitaries) \\
Symmetry generator & Killing vector $\xi$ & Self-adjoint operator $\hat{A}$ \\
What it generates & Isometric flow $\phi_t^\xi$ & Unitary group $e^{-i\hat{A}t/\hbar}$ \\
Conservation & $dQ_\xi/dt = 0$ (charge const.) & $d\langle\hat{A}\rangle/dt = 0$ (expect.\ const.) \\
Algebraic condition & $\mathcal{L}_\xi g = 0$ (Killing eq.) & $[\hat{A}, H] = 0$ (commutator) \\
\bottomrule
\end{tabular}
\end{center}
The bridge is Stone's theorem: a self-adjoint generator $\hat{A}$ produces a unitary group (norm preservation), just as a Killing vector $\xi$ produces an isometric flow ($g$-preservation). In both cases: \emph{a generator that commutes with the dynamics produces a conserved quantity}.
\end{theorem}

\begin{proof}
\textbf{Classical side.} If $\xi$ is Killing ($\mathcal{L}_\xi g = 0$) and $\nabla_\mu T^{\mu\nu} = 0$, then $Q_\xi = \int T^{\mu\nu}\xi_\nu\,\dd\Sigma_\mu$ is constant in time (Paper~2, Proposition~2.4).

\textbf{Quantum side.} If $[\hat{A}, H] = 0$ (i.e., $\hat{A}$ commutes with the Hamiltonian), then $\frac{\dd}{\dd t}\langle\psi(t)|\hat{A}|\psi(t)\rangle = \frac{i}{\hbar}\langle\psi|[H, \hat{A}]|\psi\rangle = 0$. The expectation value is conserved.

The structural parallel: in both cases, the conservation law arises because the generator commutes with the dynamics ($\mathcal{L}_\xi g = 0$ classically; $[\hat{A}, H] = 0$ quantum-mechanically). Stone's theorem provides the precise dictionary: self-adjoint operator $\leftrightarrow$ unitary group $\leftrightarrow$ conserved expectation value.
\end{proof}

\begin{remark}\label{rmk:not-analogy}
We must be precise about the scope of this correspondence. Unitarity ($U^\dagger U = 1$, norm preservation) is the quantum analog of \emph{symplecticity} (the classical flow preserves the symplectic form), not directly of Noether conservation (a specific charge is constant). The Noether analog in QM is the statement $[\hat{A}, H] = 0 \Rightarrow d\langle\hat{A}\rangle/dt = 0$, which is a \emph{consequence} of unitarity, not the same thing. The correspondence is: unitarity provides the framework within which quantum Noether conservation operates, just as symplecticity provides the framework for classical Noether conservation. The two are \emph{structurally analogous}, not identical.
\end{remark}


% ============================================================================
\section{The Non-Relativistic Case: Schr\"odinger from Bargmann}
\label{sec:schrodinger}
% ============================================================================

\subsection{The Bargmann group}

\begin{definition}[Galilei group]\label{def:galilei}
The Galilei group $\mathcal{G}_{\mathrm{Gal}}$ consists of transformations $(\tau, \vec{a}, \vec{v}, R)$ acting on spacetime by:
$t \mapsto t + \tau$, $\vec{x} \mapsto R\vec{x} + \vec{v}t + \vec{a}$,
where $\tau \in \mathbb{R}$ (time translation), $\vec{a} \in \mathbb{R}^3$ (spatial translation), $\vec{v} \in \mathbb{R}^3$ (boost), $R \in \SO(3)$ (rotation). Dimension: $1 + 3 + 3 + 3 = 10$.
\end{definition}

\begin{definition}[Bargmann group --- central extension~\cite{bargmann1954}]\label{def:bargmann}
The Galilei group admits a non-trivial \textbf{central extension} by $\mathbb{R}$:
\begin{equation}
1 \to \mathbb{R} \to \widetilde{\mathcal{G}} \to \mathcal{G}_{\mathrm{Gal}} \to 1.
\end{equation}
The extended (Bargmann) group $\widetilde{\mathcal{G}}$ has elements $(\theta, \tau, \vec{a}, \vec{v}, R)$ with composition:
\begin{equation}
(\theta_2, g_2) \cdot (\theta_1, g_1) = \left(\theta_2 + \theta_1 + \frac{m}{2\hbar}\omega(g_2, g_1),\; g_2 g_1\right)
\end{equation}
where $\omega(g_2, g_1) = \vec{v}_2 \cdot R_2\vec{a}_1 - \vec{v}_2 \cdot \vec{v}_1\tau_1$ is Bargmann's 2-cocycle.
\end{definition}

\begin{theorem}[Bargmann algebra]\label{thm:bargmann-algebra}
The Lie algebra of $\widetilde{\mathcal{G}}$ has generators $H$ (energy), $P_i$ (momentum), $K_i$ (boost), $L_{ij}$ (angular momentum), and $M$ (mass --- central). The non-vanishing brackets include:
\begin{align}
[L_{ij}, P_k] &= \delta_{ik}P_j - \delta_{jk}P_i, \label{eq:LP-gal}\\
[L_{ij}, K_k] &= \delta_{ik}K_j - \delta_{jk}K_i, \label{eq:LK-gal}\\
[K_i, H] &= P_i, \label{eq:KH}\\
[K_i, P_j] &= \delta_{ij}\,M, \label{eq:KP-gal}
\end{align}
where $M$ is central: $[M, \cdot\,] = 0$. The bracket~\eqref{eq:KP-gal} is the defining relation of the central extension. It has no classical analog (for the unextended Galilei group, $[K_i, P_j] = 0$).
\end{theorem}

\begin{remark}[Mass as central charge]\label{rmk:mass-central}
In Paper~1, the speed of light $c$ is the structural constant of the Poincar\'e kernel. Here, the mass $m$ enters as the \emph{central charge} of the Bargmann extension. In both cases, a fundamental physical quantity (speed, mass) is a \emph{parameter of the symmetry algebra}, not a property of a particular object.
\end{remark}

\subsection{The three axioms for quantum mechanics}

\begin{axiom}[Bargmann equivariance]\label{ax:Q1}
Time evolution is implemented by a unitary representation $U: \widetilde{\mathcal{G}} \to \mathcal{U}(\mathcal{H})$ of the Bargmann group on a Hilbert space $\mathcal{H}$.
\end{axiom}

\begin{axiom}[Unitarity]\label{ax:Q2}
The time evolution $U(t) = e^{-iHt/\hbar}$ preserves the norm: $\|U(t)\psi\| = \|\psi\|$ for all $t$ and $\psi \in \mathcal{H}$. Equivalently, $H$ is self-adjoint.
\end{axiom}

\begin{axiom}[Irreducibility]\label{ax:Q3}
The representation $U$ is irreducible: no proper closed subspace of $\mathcal{H}$ is invariant under all $U(g)$.
\end{axiom}

\begin{remark}
These axioms parallel the GR axioms of Paper~2: (Q1) $\leftrightarrow$ (A1) [group equivariance]; (Q2) $\leftrightarrow$ (A2) [conservation]; (Q3) $\leftrightarrow$ (A3) [a simplicity constraint selecting the minimal structure].
\end{remark}

\subsection{Uniqueness via Casimirs and Schur's lemma}

\begin{theorem}[Schur's lemma]\label{thm:schur}
If $\pi: G \to \mathcal{U}(\mathcal{H})$ is an irreducible unitary representation and $T: \mathcal{H} \to \mathcal{H}$ commutes with all $\pi(g)$ ($[T, \pi(g)] = 0$ for all $g \in G$), then $T = \lambda \cdot \mathrm{Id}$ for some scalar $\lambda$.
\end{theorem}

\begin{definition}[Casimir operators]\label{def:casimirs-bargmann}
The \textbf{Casimir operators} of the Bargmann algebra are the elements of the center of the universal enveloping algebra $U(\widetilde{\mathfrak{g}})$. They commute with all generators and hence, on any representation, commute with all $U(g)$. For the Bargmann algebra, the Casimirs are:
\begin{align}
C_1 &= M, \label{eq:C1-gal}\\
C_2 &= H - \frac{P^2}{2M}, \label{eq:C2-gal}\\
C_3 &= \left(\vec{L} - \frac{\vec{K} \times \vec{P}}{M}\right)^2. \label{eq:C3-gal}
\end{align}
Here $C_1$ is the mass, $C_2$ is the internal (non-translational) energy, and $C_3$ encodes the spin.
\end{definition}

\begin{theorem}[Schr\"odinger equation from Bargmann representations]\label{thm:schrodinger}
Axioms~Q1--Q3 uniquely determine the free Schr\"odinger equation:
\begin{equation}\label{eq:schrodinger}
\boxed{i\hbar\frac{\partial\psi}{\partial t} = -\frac{\hbar^2}{2m}\nabla^2\psi + E_0\psi}
\end{equation}
where $E_0$ is the internal energy (a representation-dependent constant). Adding an external potential $V(\vec{x})$ (which breaks Galilean symmetry locally) gives the full Schr\"odinger equation $i\hbar\partial_t\psi = (-\hbar^2\nabla^2/2m + V)\psi$.
\end{theorem}

\begin{proof}
The proof proceeds in four steps, mirroring the structure of the GR derivation (Paper~2).

\textbf{Step 1} (Casimir eigenvalues from Schur). By Axiom~Q3 (irreducibility), Schur's lemma (Theorem~\ref{thm:schur}) applies to each Casimir:
\begin{equation}
C_1 = m \cdot \mathrm{Id}, \qquad C_2 = E_0 \cdot \mathrm{Id}, \qquad C_3 = s(s+1)\hbar^2 \cdot \mathrm{Id},
\end{equation}
where $m > 0$ is the mass, $E_0 \in \mathbb{R}$ is the internal energy, and $s \in \{0, \frac{1}{2}, 1, \ldots\}$ is the spin.

\textbf{Step 2} (Hamiltonian from $C_2$). The Casimir equation $C_2 = E_0 \cdot \mathrm{Id}$ reads:
\begin{equation}
H - \frac{P^2}{2m} = E_0 \cdot \mathrm{Id}, \qquad\text{hence}\qquad H = \frac{P^2}{2m} + E_0.
\end{equation}
The Hamiltonian is \emph{determined} by the representation theory (the Casimir eigenvalue), not postulated.

\textbf{Step 3} (Position representation). By Axiom~Q2 (unitarity) and the canonical commutation relations $[\hat{x}_i, P_j] = i\hbar\delta_{ij}$ (which follow from the Bargmann algebra via exponentiation~\cite{bargmann1954}): the Schr\"odinger representation is $P_i = -i\hbar\partial/\partial x_i$ on $\mathcal{H} = L^2(\mathbb{R}^3)$. The Hamiltonian becomes:
\begin{equation}
H = \frac{P^2}{2m} + E_0 = -\frac{\hbar^2}{2m}\nabla^2 + E_0.
\end{equation}

\textbf{Step 4} (Time evolution). By Stone's theorem (Theorem~\ref{thm:stone}), the unitary evolution $U(t) = e^{-iHt/\hbar}$ gives $i\hbar\partial_t\psi = H\psi$, which is the Schr\"odinger equation~\eqref{eq:schrodinger}.
\end{proof}

\begin{remark}[What was not used]\label{rmk:no-lagrangian-qm}
No action was minimized. No classical Hamiltonian was ``quantized.'' The Schr\"odinger equation emerged from: (1)~the Bargmann group (symmetry), (2)~unitarity (conservation), (3)~Schur's lemma (uniqueness). The Casimir $C_2 = H - P^2/2m$ plays the role of the Bianchi identity in GR: it is the orbital constraint that determines the dynamics.

We acknowledge that the derivation of the Schr\"odinger equation from the Bargmann group is a well-known result (Bargmann~\cite{bargmann1954}, L\'evy-Leblond~\cite{levyleblond1967}). The novelty of this paper lies not in the derivation itself but in: (a)~placing it in the three-step orbital pattern alongside GR (Paper~2); (b)~the unitarity--Noether correspondence (Theorem~\ref{thm:correspondence}); (c)~the Bianchi--Casimir correspondence (Theorem~\ref{thm:bianchi-casimir}).
\end{remark}

\begin{remark}[Stone--von Neumann uniqueness]\label{rmk:stone-vn}
The position representation used in Step~3 (i.e., $P_i = -i\hbar\partial/\partial x_i$ on $L^2(\mathbb{R}^3)$) is not merely one possible choice: the \textbf{Stone--von Neumann theorem}~\cite{stonevn1931,vonneumann1931} guarantees that it is the \emph{unique} (up to unitary equivalence) irreducible representation of the canonical commutation relations $[\hat{x}_i, P_j] = i\hbar\delta_{ij}$ on a separable Hilbert space (in the Weyl exponentiated form). This provides a second uniqueness result, complementing Schur's lemma: Schur determines the Hamiltonian ($H = P^2/2m + E_0$); Stone--von Neumann determines the representation space ($L^2(\mathbb{R}^3)$). Together, they completely fix the Schr\"odinger equation.
\end{remark}

\begin{remark}[The Bargmann cocycle is the Galilei-boost phase]\label{rmk:bargmann-phase}
The non-triviality of the central extension (Definition~\ref{def:bargmann}) is directly observable. Under a
Galilei boost $\vec x'=\vec x-\vec v\,t$, $t'=t$, the free Schr\"odinger equation
$i\hbar\partial_t\psi=-\tfrac{\hbar^2}{2m}\nabla^2\psi$ is covariant only if the wavefunction transforms with a
phase (in the standard convention)
\[
\psi'(\vec x',t')=\exp\!\Big(\tfrac{i}{\hbar}\big(m\,\vec v\cdot\vec x-\tfrac12 m v^2 t\big)\Big)\,\psi(\vec x,t).
\]
Direct substitution confirms covariance, and the phase cannot be gauged away. This phase \emph{is} Bargmann's
$2$-cocycle $\tfrac{m}{2\hbar}\omega$ of Definition~\ref{def:bargmann}, with the mass $m$ as central charge:
composing two boosts, the phases combine with the cocycle discrepancy $\propto m\,\omega(g_2,g_1)$, which is
precisely the obstruction to the Galilei group acting by a genuine (non-projective) representation. The
unextended Galilei group ($[K_i,P_j]=0$) carries no such phase and hence no nontrivial mass; the central
extension is therefore not a technical convenience but the algebraic home of mass.
\end{remark}

\begin{remark}[The potential $V$]
The external potential $V(\vec{x})$ breaks the translational symmetry of the Galilei group (just as the stress-energy tensor $T_{\mu\nu}$ breaks the Poincar\'e symmetry in GR). In both cases, the symmetry-breaking source is an input to the equation, not derived from the framework.
\end{remark}


% ============================================================================
\section{The Relativistic Case: Wigner Classification}
\label{sec:wigner}
% ============================================================================

For relativistic quantum mechanics, the Bargmann group is replaced by the Poincar\'e group, and the Casimirs change accordingly.

\begin{theorem}[Poincar\'e Casimirs]\label{thm:poincare-casimirs}
The irreducible unitary representations of the Poincar\'e group are classified by two Casimir invariants:
\begin{align}
C_1 &= P^\mu P_\mu = -m^2c^2, \label{eq:C1-poinc}\\
C_2 &= W^\mu W_\mu = -m^2c^2\,s(s+1)\hbar^2, \label{eq:C2-poinc}
\end{align}
where $W^\mu = \frac{1}{2}\epsilon^{\mu\nu\rho\sigma}M_{\nu\rho}P_\sigma$ is the Pauli--Lubanski pseudovector. The label $(m, s)$ specifies the mass and spin.
\end{theorem}

\begin{theorem}[Wave equations from Wigner classification]\label{thm:wigner-eqs}
On an irreducible representation with mass $m$ and spin $s$, the Casimir eigenvalue $C_1 = -m^2c^2$ determines the wave equation. In the position representation ($P_\mu = -i\hbar\partial_\mu$):
\begin{equation}
C_1\,\psi = -m^2c^2\,\psi \qquad\Longrightarrow\qquad \left(\Box - \frac{m^2c^2}{\hbar^2}\right)\psi = 0,
\end{equation}
where $\Box = \eta^{\mu\nu}\partial_\mu\partial_\nu = -c^{-2}\partial_t^2 + \nabla^2$ is the d'Alembertian. This is the \textbf{Klein--Gordon equation}.

For different spins, the little group (stabilizer of a reference momentum) determines the field content:
\begin{center}\small
\begin{tabular}{ccll}\toprule
$m$ & $s$ & \textbf{Little group} & \textbf{Wave equation} \\\midrule
$> 0$ & $0$ & $\SO(3)$ & Klein--Gordon: $(\Box - m^2c^2/\hbar^2)\phi = 0$ \\
$> 0$ & $\frac{1}{2}$ & $\SU(2)$ & Dirac: $(i\hbar\gamma^\mu\partial_\mu - mc)\psi = 0$ \\
$= 0$ & $1$ & $\ISO(2)$ & Maxwell: $\partial_\mu F^{\mu\nu} = 0$ \\
$= 0$ & $2$ & $\ISO(2)$ & Linearized Einstein \\
\bottomrule
\end{tabular}
\end{center}
\end{theorem}

\begin{proof}[Proof for Klein--Gordon]
The quadratic Casimir on the representation: $C_1 = P^\mu P_\mu = (-i\hbar\partial^\mu)(-i\hbar\partial_\mu) = -\hbar^2\Box$. By Schur's lemma on the irreducible representation: $C_1 = -m^2c^2 \cdot \mathrm{Id}$. Hence $-\hbar^2\Box\psi = -m^2c^2\psi$, i.e.\ $\Box\psi = (m^2c^2/\hbar^2)\psi$, giving $(\Box - m^2c^2/\hbar^2)\psi = 0$.
\end{proof}

\begin{proof}[Proof for Dirac]
The Dirac equation arises from factorizing the Klein--Gordon operator. Taking $\gamma$-matrices with the Clifford algebra $\{\gamma^\mu, \gamma^\nu\} = -2\eta^{\mu\nu}$ (so that $(i\hbar\gamma^\mu\partial_\mu)^2 = \hbar^2\Box$), one has $(i\hbar\gamma^\mu\partial_\mu + mc)(i\hbar\gamma^\nu\partial_\nu - mc) = \hbar^2\!\left(\Box - \tfrac{m^2c^2}{\hbar^2}\right)$. Hence the first-order equation $(i\hbar\gamma^\mu\partial_\mu - mc)\psi = 0$ --- the Dirac equation --- implies Klein--Gordon, and its solutions form a spin-$\frac{1}{2}$ representation of the Poincar\'e group~\cite{wigner1939,levyleblond1967}.
\end{proof}

\begin{remark}[Spin without relativity]
L\'evy-Leblond~\cite{levyleblond1967} showed that the same factorization procedure applied to the Bargmann Casimir $C_2 = H - P^2/2m$ produces a four-component equation (the L\'evy-Leblond equation) whose solutions carry spin-$\frac{1}{2}$ and yield the correct gyromagnetic ratio $g = 2$. Spin is a consequence of the \emph{rotation subgroup} $\SO(3)$ (common to both Galilei and Poincar\'e), not of relativity.
\end{remark}


% ============================================================================
\section{The Bianchi--Casimir Correspondence}
\label{sec:correspondence}
% ============================================================================

This section contains the central structural result of the paper.

\begin{theorem}[Bianchi--Casimir correspondence]\label{thm:bianchi-casimir}
The Bianchi identity (Paper~2) and the Casimir eigenvalue equation (this paper) play structurally identical roles in their respective theories:

\begin{center}\small
\begin{tabular}{lll}\toprule
\textbf{Role} & \textbf{GR (Paper~2)} & \textbf{QM (this paper)} \\\midrule
Symmetry group & $\Diff(\mathcal{M})$ & Bargmann / Poincar\'e \\
State space & Lorentzian metrics $g$ & Hilbert space $\mathcal{H}$ \\
Equivariance & Natural concomitant & Unitary representation \\
Conservation & $\nabla_\mu T^{\mu\nu} = 0$ & $\|U\psi\| = \|\psi\|$ \\
\textbf{Orbital constraint} & $\nabla_\mu G^{\mu\nu} \equiv 0$ \textbf{(Bianchi)} & $C_i = \lambda_i \cdot \mathrm{Id}$ \textbf{(Casimir)} \\
\textbf{Uniqueness theorem} & Curiel/Navarro-Sancho & Schur's lemma \\
Output & $G_{\mu\nu} + \Lambda g_{\mu\nu} = 8\pi T_{\mu\nu}$ & $i\hbar\partial_t\psi = H\psi$ \\
\bottomrule
\end{tabular}
\end{center}

Both constraints are \textbf{orbital integrability conditions}: they specify which orbit of the symmetry group the physical system occupies, and the uniqueness theorem converts this orbit specification into a unique equation.
\end{theorem}

\begin{proof}
The proof establishes the structural parallel step by step.

\textbf{Step 1 (Equivariance $\to$ constraint).} In GR: $\Diff$-equivariance (Axiom~A1) requires $\mathcal{E}_{\mu\nu}$ to be a natural concomitant. The Bianchi identity $\nabla_\mu G^{\mu\nu} \equiv 0$ is a geometric identity constraining which natural concomitants are divergence-free. In QM: Bargmann equivariance (Axiom~Q1) requires $U$ to be a unitary representation. The Casimir operators $C_i$ commute with all $U(g)$ by construction, constraining which representations are irreducible.

\textbf{Step 2 (Constraint $\to$ uniqueness).} In GR: the constraint (divergence-free + second-order + natural) determines $\mathcal{E} = aG + bg$ by Curiel's theorem. In QM: the constraint ($C_i = \lambda_i$ on an irreducible representation) determines $H = P^2/2m + E_0$ by Schur's lemma.

\textbf{Step 3 (Uniqueness $\to$ equation).} In GR: $\mathcal{E} = aG + bg = \kappa T$ gives the Einstein equation. In QM: $H = P^2/2m + E_0$ gives the Schr\"odinger equation via Stone's theorem.

\textbf{Step 4 (Conservation is automatic).} In GR: $\nabla_\mu(aG + bg)^{\mu\nu} = 0$ (Bianchi + metric compatibility) guarantees $\nabla_\mu T^{\mu\nu} = 0$. In QM: $H = H^\dagger$ (self-adjointness, forced by the Casimir being real) guarantees unitarity $U(t)^\dagger U(t) = 1$.
\end{proof}

\begin{remark}[Depth of the correspondence]
The parallel is not superficial. In both cases:
\begin{enumerate}[label=(\roman*)]
    \item The constraint is an \emph{identity} --- it holds for all metrics (Bianchi) or for all irreducible representations (Casimir), not just specific solutions.
    \item The constraint reduces the space of possible equations from infinite-dimensional to finite-dimensional (the space of constants $a, b$ for GR; the space of labels $m, s, E_0$ for QM).
    \item The constraint is \emph{forced by the group structure}, not postulated. The Bianchi identity follows from the Jacobi identity for $\nabla$; the Casimir eigenvalue follows from the definition of irreducibility.
    \item Conservation is a \emph{consequence}, not a premise: it follows automatically from the equation determined by the constraint.
\end{enumerate}
\end{remark}


% ============================================================================
\section{The Classification Grid}
\label{sec:grid}
% ============================================================================

\begin{theorem}[Two-parameter classification]\label{thm:grid}
The fundamental theories of physics are classified by the pair $(\kappa, \rho)$:
\begin{equation}\label{eq:grid}
\begin{array}{c|cc}
& \kappa = 10 & \kappa < 10 \\ \hline
\text{Regime } (\alpha) & \textbf{SR} & \textbf{GR} \\[3pt]
\text{Regime } (\beta) & \textbf{QM} & \textbf{Quantum Gravity}
\end{array}
\end{equation}
Each cell is determined by the three-step orbital pattern with the appropriate group, conservation principle, and uniqueness theorem. The transitions between cells are:
\begin{itemize}
    \item SR $\to$ GR: reduce $\kappa$ from $10$ to $< 10$ (introduce curvature).
    \item SR $\to$ QM: change regime $\alpha \to \beta$ (deterministic $\to$ stochastic).
    \item QM $\to$ QG: reduce $\kappa$ (quantum mechanics on curved spacetime).
    \item GR $\to$ QG: change regime $\alpha \to \beta$ (quantize gravity).
\end{itemize}
\end{theorem}

\begin{remark}[Limits]
QM $\to$ SR when $\hbar \to 0$ (the stochastic coupling $P(\cdot|g)$ concentrates on the deterministic value). QG $\to$ GR when $\hbar \to 0$. QG $\to$ QM when $\kappa \to 10$ (flat limit). The fourth corner (QG) is constrained by requiring consistency with all three known corners.
\end{remark}


% ============================================================================
\section{Discussion}
\label{sec:discussion}
% ============================================================================

\subsection{What the framework does not do}

\begin{enumerate}[label=(\roman*)]
    \item \textbf{$\hbar$ is not derived.} Like $c$ (Paper~1) and $G, \Lambda$ (Paper~2), $\hbar$ is a structural constant of the symmetry algebra (the central extension parameter). Its value is an empirical input.
    \item \textbf{The measurement problem is not solved.} The framework provides a geometric language (measurement $=$ projection onto orbits of the observable's symmetry group), but not a mechanism for wave function collapse. It is compatible with any interpretation (Copenhagen, many-worlds, etc.).
    \item \textbf{Quantum field theory requires a separate treatment.} The multi-particle extension (Fock space, spin-statistics, CPT) requires the cluster decomposition principle, which goes beyond single-particle QM. This is addressed in Paper~3b~\cite{nembe2026p3b}.
\end{enumerate}

\subsection{The three-step pattern across the series}

\begin{center}\small
\begin{tabular}{lccc}\toprule
\textbf{Step} & \textbf{Paper 1 (SR)} & \textbf{Paper 2 (GR)} & \textbf{Paper 3a (QM)} \\\midrule
1. Group & Poincar\'e & $\Diff(\mathcal{M})$ & Bargmann / Poincar\'e \\
2. Conservation & $dQ/dt = 0$ & $\nabla_\mu T^{\mu\nu} = 0$ & $\|U\psi\| = \|\psi\|$ \\
3. Uniqueness & Max-sym classif. & Curiel/N-S & Schur's lemma \\
\midrule
Orbital constraint & Mass-shell & Bianchi & Casimir \\
Output & Lorentz + $E^2 = \ldots$ & $G + \Lambda g = 8\pi T$ & $i\hbar\partial_t\psi = H\psi$ \\
Lagrangian & Not used & Not used & Not used \\
\bottomrule
\end{tabular}
\end{center}

\subsection{Connection to Paper 1}

The mass-shell $E^2 = p^2c^2 + m^2c^4$ (Paper~1, Theorem~5.3, from the Poincar\'e coadjoint orbit) and the Klein--Gordon equation $(\Box - m^2c^2/\hbar^2)\phi = 0$ (this paper, Theorem~\ref{thm:wigner-eqs}) arise from the \emph{same} Casimir $C_1 = P^\mu P_\mu = -m^2c^2$. The difference: in regime $(\alpha)$, the Casimir gives an orbit constraint on a \emph{point} (the 4-momentum lives on the mass hyperboloid); in regime $(\beta)$, it gives an operator equation on a \emph{wave function} (the field satisfies the Klein--Gordon equation). The group structure is identical; the representation space changes.


% ============================================================================
\section{Conclusion}
\label{sec:conclusion}
% ============================================================================

We have derived the Schr\"odinger, Klein--Gordon, and Dirac equations from three axioms (Bargmann/Poincar\'e equivariance, unitarity, irreducibility) without any variational principle, following the same three-step orbital pattern as Papers~1 (SR) and~2 (GR).

The central results are: (i)~the unitarity--Noether correspondence (Theorem~\ref{thm:correspondence}), showing that quantum unitarity is the structural analog of classical Noether conservation; (ii)~the Bianchi--Casimir correspondence (Theorem~\ref{thm:bianchi-casimir}), showing that the Bianchi identity (GR) and the Casimir eigenvalue (QM) play structurally identical roles as orbital integrability conditions; (iii)~the $(\kappa, \rho)$ classification grid (Theorem~\ref{thm:grid}), placing SR, GR, QM, and quantum gravity in a unified framework.

The pattern is now established for three of the four corners of the grid. The fourth corner --- quantum gravity --- is the subject of Paper~4~\cite{nembe2026p4}, which will formalize the universal pattern as a metatheorem and derive constraints on the missing theory.


% ============================================================================
\section*{Acknowledgments}
% ============================================================================

The author thanks the Roots Insights Laboratory and the National Institute of Management Sciences for support.


% ============================================================================
\begin{thebibliography}{30}

\bibitem{schrodinger1926} E.~Schr\"odinger, ``Quantisierung als Eigenwertproblem,'' \emph{Ann.\ Phys.}\ \textbf{384}, 361--376 (1926).

\bibitem{dirac1928} P.A.M.~Dirac, ``The quantum theory of the electron,'' \emph{Proc.\ R.\ Soc.\ A}\ \textbf{117}, 610--624 (1928).

\bibitem{feynman1948} R.P.~Feynman, ``Space-time approach to non-relativistic quantum mechanics,'' \emph{Rev.\ Mod.\ Phys.}\ \textbf{20}, 367--387 (1948).

\bibitem{wigner1939} E.P.~Wigner, ``On unitary representations of the inhomogeneous Lorentz group,'' \emph{Ann.\ Math.}\ \textbf{40}, 149--204 (1939).

\bibitem{bargmann1954} V.~Bargmann, ``On unitary ray representations of continuous groups,'' \emph{Ann.\ Math.}\ \textbf{59}, 1--46 (1954).

\bibitem{levyleblond1967} J.-M.~L\'evy-Leblond, ``Nonrelativistic particles and wave equations,'' \emph{Comm.\ Math.\ Phys.}\ \textbf{6}, 286--311 (1967).

\bibitem{musielak2021} Z.E.~Musielak, ``Nonrelativistic fundamental quantum and classical wave equations,'' \emph{Int.\ J.\ Mod.\ Phys.\ A}\ \textbf{36}, 2150042 (2021).

\bibitem{stone1932} M.H.~Stone, ``On one-parameter unitary groups in Hilbert space,'' \emph{Ann.\ Math.}\ \textbf{33}, 643--648 (1932).

\bibitem{reed1972} M.~Reed and B.~Simon, \emph{Methods of Modern Mathematical Physics}, Vol.~I (Academic Press, 1972).

\bibitem{mackey1968} G.W.~Mackey, \emph{Induced Representations of Groups and Quantum Mechanics} (Benjamin, 1968).

\bibitem{inonu1953} E.~In\"on\"u and E.P.~Wigner, ``On the contraction of groups and their representations,'' \emph{Proc.\ Nat.\ Acad.\ Sci.}\ \textbf{39}, 510--524 (1953).

\bibitem{weinberg1995} S.~Weinberg, \emph{The Quantum Theory of Fields}, Vol.~1 (Cambridge Univ.\ Press, 1995).

\bibitem{jauch1968} J.M.~Jauch, \emph{Foundations of Quantum Mechanics} (Addison-Wesley, 1968).

\bibitem{nembe2026p1} J.~Nemb\'e, ``Special relativity from orbital structure'' (Paper~1, 2026).

\bibitem{nembe2026p2} J.~Nemb\'e, ``General relativity from orbital structure'' (Paper~2, 2026).

\bibitem{nembe2026p3b} J.~Nemb\'e, ``Quantum field theory from orbital structure'' (Paper~3b, in preparation).

\bibitem{nembe2026p4} J.~Nemb\'e, ``The equations of physics from orbital structure: a metatheorem'' (Paper~4, in preparation).

\bibitem{nembe2026vol17} J.~Nemb\'e, ``Twin General Relativity,'' \emph{ROOTS-INSIGHTS} Vol.~17 (2026).

\bibitem{stonevn1931} M.H.~Stone, ``Linear transformations in Hilbert space. III.\ Operational methods and group theory,'' \emph{Proc.\ Nat.\ Acad.\ Sci.}\ \textbf{16}, 172--175 (1930).

\bibitem{vonneumann1931} J.~von Neumann, ``Die Eindeutigkeit der Schr\"odingerschen Operatoren,'' \emph{Math.\ Ann.}\ \textbf{104}, 570--578 (1931).

\bibitem{strocchi2008} F.~Strocchi, \emph{An Introduction to the Mathematical Structure of Quantum Mechanics}, 2nd ed.\ (World Scientific, 2008).

\bibitem{giulini2009} D.~Giulini, ``On Galilei invariance in quantum mechanics and the Bargmann superselection rule,'' \emph{Ann.\ Phys.}\ \textbf{316}, 393--428 (2005).

\end{thebibliography}

\end{document}
